% Part: first-order-logic
% Chapter: tableaux
% Section: soundness-identity

\documentclass[../../../include/open-logic-section]{subfiles}

\begin{document}

\olfileid[fr]{fol}{tab}{sid}

\olsection{Correction avec \usetoken{S}{identity}}

\begin{prop}
Les !!^{tableau}s munis des règles de l'identité sont corrects : aucun
!!{tableau} fermé n'est satisfaisable.
\end{prop}

\begin{proof}
Comme précédemment, il suffit de montrer que, si !!a{tableau} possède
une branche satisfaisable, la branche obtenue en appliquant à celle-ci
l'une des règles de~$\eq$ est encore satisfaisable. Soit $\Gamma$
l'ensemble des !!{signed formula}s de cette branche, et soit $\Struct{M}$
!!a{structure} qui satisfait~$\Gamma$.

Supposons que la branche soit prolongée au moyen de~$\eq$, c'est-à-dire
par l'ajout de la !!{formula} signée
$\sFmla{\True}{\eq[t][t]}$. On a trivialement
$\Sat{M}{\eq[t][t]}$; par conséquent, $\Struct{M}$ satisfait aussi
$\Gamma \cup \{\sFmla{\True}{\eq[t][t]}\}$.

Si la branche est prolongée au moyen de $\TRule{\True}{\eq}$, nous
ajoutons la !!{formula} signée $\sFmla{\True}{!A(t_2)}$, tandis que
$\Gamma$ contient déjà $\sFmla{\True}{\eq[t_1][t_2]}$ et
$\sFmla{\True}{!A(t_1)}$. Nous avons donc
$\Sat{M}{\eq[t_1][t_2]}$ et $\Sat{M}{!A(t_1)}$. Soit $s$ une
assignation de variables telle que $s(x) = \Value{t_1}{M}$. D'après
\olref[syn][ass]{prop:sentence-sat-true}, $\Sat{M}{!A(t_1)}[s]$.
Comme $\varAssign{s}{s}{x}$, la
\olref[syn][ext]{prop:ext-formulas} donne $\Sat{M}{!A(x)}[s]$.
Puisque $\Sat{M}{\eq[t_1][t_2]}$, on a
$\Value{t_1}{M} = \Value{t_2}{M}$, et donc
$s(x) = \Value{t_2}{M}$. Une nouvelle application de
\olref[syn][ext]{prop:ext-formulas} donne alors
$\Sat{M}{!A(t_2)}[s]$. Enfin,
\olref[syn][ass]{prop:sentence-sat-true} entraîne
$\Sat{M}{!A(t_2)}$. Le cas de $\TRule{\False}{\eq}$ se traite de la
même manière.
\end{proof}

\begin{editorial}
Dans le cas de $\TRule{\True}{\eq}$, la source donne le signe
générique~$S$ à la conclusion. Le signe~$\True$ exigé par la règle est
rétabli; le cas du signe~$\False$ reste traité séparément par analogie.
\end{editorial}

\end{document}
